\documentclass[11pt,letterpaper]{article}
\usepackage[margin=0.88in,headheight=15pt]{geometry}
\usepackage[T1]{fontenc}
\usepackage{lmodern,microtype,amsmath,amssymb,amsthm,mathtools}
\usepackage{booktabs,tabularx,array,xcolor,fancyhdr}
\usepackage[colorlinks=true,linkcolor=blue!45!black,urlcolor=blue!50!black,citecolor=blue!45!black]{hyperref}
\definecolor{ink}{RGB}{28,49,68}
\pagestyle{fancy}\fancyhf{}
\setlength{\parindent}{0pt}\setlength{\parskip}{5pt}
\setlength{\emergencystretch}{2em}
\setcounter{tocdepth}{1}\numberwithin{equation}{section}
\theoremstyle{plain}
\newtheorem{theorem}{Theorem}[section]
\newtheorem{lemma}[theorem]{Lemma}
\newtheorem{proposition}[theorem]{Proposition}
\theoremstyle{remark}\newtheorem{remark}[theorem]{Scope}
\DeclareMathOperator{\Cov}{Cov}\DeclareMathOperator{\supp}{supp}
\DeclareMathOperator{\Part}{Part}\DeclareMathOperator{\diam}{diam}
\DeclareMathOperator{\Ad}{Ad}\DeclareMathOperator{\dist}{dist}
\newcommand{\dd}{\,\mathrm d}\newcommand{\eps}{\varepsilon}
\newcommand{\norm}[1]{\left\lVert#1\right\rVert}
\newcommand{\abs}[1]{\left|#1\right|}
\newcommand{\R}{\mathbb R}\newcommand{\E}{\mathbb E}
\title{Uniform averaging rows and finite-jet response}\author{}\date{}
\providecommand{\mathscr}{\mathcal}
\begin{document}\maketitle
\section{The exact target and its saved three rows}

Equation names LM, MT, AE8, AE10, CM, CB, and NM12 refer to~\cite{W}: CM on pp.~326--327, CB on pp.~666--672, MT on pp.~686--698, LM on pp.~773--784, AE8 on pp.~932--935, AE10 on pp.~1070--1072 and NM12 on pp.~1216--1218. The note~\cite[\S3]{CE} supplies the raw connected
estimate but leaves the incoming bound open. For a fixed original history
$\omega$ and all its simultaneous walls $s$, set, wherever the actual maps
are admitted,
\begin{equation}
 W_{\omega,s}=P_{\rm old}\bigl(q_{\rm old}(\Gamma^R_{\omega,s})\bigr),
 \qquad
 F_{j,k;\omega,s}
 =R_j\bigl(W_{\omega,s},J_{\rm phys}(W_{\omega,s})\bigr).
 \label{eq:actual}
\end{equation}
Here $q_{\rm old}$ includes every determining row, prescribed prefix,
junction, endpoint frame, and literal boundary datum. $P_{\rm old}$ is the
original constrained Wilson minimum with that complete record. The separately
born proper differences and lower/contact packet remain part of the state.
LM46 gives the physical endpoint identity on its common domain. It does not
give a uniform norm on the whole wall polydisc.

Write $r=L^{j-k}$ and $\rho=g_k/g_j$. For a clean physical radius $\alpha$,
the saved canonical weights in MT44 are
\begin{equation}
\begin{array}{c|cc}
 & d=2&d=3\\ \hline
\text{quadratic}&r^s\rho^2&r^s\rho^2\\[2pt]
\text{cubic}&r^\beta(g_k/\alpha)\rho^2&r^\beta\rho^3\\[2pt]
\text{quartic tail}&r^2(g_k/\alpha)^2\rho^2&r^2(g_k/\alpha)\rho^3
\end{array}
\label{eq:weights}
\end{equation}
The root count and dimension-six normalization have already been included.
For an R-origin coefficient, its recorded shrinking radii and promotion wait
replace the clean-radius simplification; they cannot be reset at entry.

The target is an actual incoming envelope with these saved factors, its
finite current jets, and the coefficient/anchor moment sums. Its proof also
needs the total quadratic/cubic covariant cancellation before absolute values
are taken. A supremum bound for a native action member does not supply those
moments. Likewise, an old derivative at fixed canonical directions does not
bound a sum over all source anchors.

The regular-window proof uses old dimensionless field and own-current sizes
of order $r^2g_kq$ and $r^3g_kq$. The immediate question below is whether
the \emph{actual old target} has those sizes. A current-scale unweighted norm
for the entire prepared microscopic field would be a stronger sufficient
condition in some directions. It is useful to avoid requiring that stronger
condition before examining the old averaging map itself.

\section{A direct estimate for complete uniform old averaging rows}
\label{sec:average}

\subsection{The actual flat differential and its domain}

The raw group mean has the outer form
\begin{equation}
 M(\{P_i\},R)
 =\exp\left\{\sum_i w_i\log(P_iR^{-1})\right\}R,
 \quad w_i\ge0,\quad\sum_iw_i=1.
 \label{eq:rawmean}
\end{equation}
Its reference word $R$ and its averaged endpoint contours are parts of the
map. The primary definition is in~\cite{sf1}, pp.\ 253--254,
(0.4), (0.8), (0.11)--(0.12). CM15--18 and CB9--12 derive its flat
differential with these endpoints retained.

Fix a complete old averaging plaquette: its four bonds and required endpoint
patches use the same uniform nested raw hierarchy. Let $a$ be the input mesh,
$r=Ba$ its old output mesh, and $B$ an odd blocking product. In physical
connection units,
\begin{equation}
 Dq_{{\rm raw},\mu}(0)=Q_0L_{\mu,B}-d_{r,\mu}\mathfrak r,
 \qquad L_{\mu,B}=\frac1B\sum_{t=0}^{B-1}\tau_{tae_\mu}.
\end{equation}
Here $\mathfrak r$ is the actual averaged endpoint connector. Taking the flat
curl cancels this exact coboundary:
\begin{equation}
 d_rDq_{\rm raw}(0)=Q_2d_a,
 \qquad Q_{2,\mu\nu}=Q_0L_{\mu,B}L_{\nu,B}.
 \label{eq:flatcurl}
\end{equation}
This is a statement about the differential of the original map. It is not
associativity of its nonlinear group means. The positive filters in
\eqref{eq:flatcurl} compose to the direct filters with $B=r/a$: at each
coordinate, the block-center digits and the line-segment digits enumerate
the corresponding direct interval exactly once.

\begin{lemma}[Old-target overlap estimate]
For a fixed two-cochain component, write
\[
 (Q_2F)(x)=\sum_yK_r(x,y)F(y).
\]
On the complete uniform row family just specified,
\begin{equation}
 K_r\ge0,\quad \sum_yK_r(x,y)=1,\quad
 K_r(x,y)\le C(a/r)^4,\quad
 \supp K_r(x,\cdot)\subset\{|x-y|\le Cr\}.
 \label{eq:kernel}
\end{equation}
Partition its input readset into actual assigned cells $\Delta$ and put
\begin{equation}
 w_{x\Delta}=\sum_{y\in\Delta}K_r(x,y),\qquad
 c_\Delta=\sigma_\Delta^2\sup_{y\in\Delta}|F(y)|.
\end{equation}
Then
\begin{equation}
 r^2|(Q_2F)(x)|
 \le\sum_\Delta w_{x\Delta}(r/\sigma_\Delta)^2c_\Delta,
 \qquad \sum_\Delta w_{x\Delta}=1.
 \label{eq:overlap}
\end{equation}
Moreover, with physical counting volume $|A|_a=a^4\#A$,
\begin{equation}
 w_{x\Delta}
 \le C\min\{1,|\Delta\cap\supp K_r(x,\cdot)|_a/r^4\}.
 \label{eq:volume}
\end{equation}
For a cell of side $\sigma\le r$ and volume at most $C\sigma^4$, its
contribution to \eqref{eq:overlap} is at most $C(\sigma/r)^2c_\Delta$.
\end{lemma}

\begin{proof}
The coefficients of $Q_0L_{\mu,B}L_{\nu,B}$ are $B^{-6}$ times
multiplicities. At fixed input and output, at most $B^2$ pairs of line
shifts contribute, giving $B^{-4}$. Positivity, row mass one, and the
support bound follow from the normalized block and line factors. Summing
these coefficients over the cell proves \eqref{eq:volume}. Insert
$|F(y)|\le\sigma_\Delta^{-2}c_\Delta$ separately on each cell to obtain
\eqref{eq:overlap}. The direct-filter composition above makes the constants
independent of the number of intermediate uniform steps.
\end{proof}

On the full periodic uniform lattice there is also the exact physical
column identity
\begin{equation}
 r^4\sum_xK_r(x,y)=a^4.
\end{equation}
It follows because each fine site has block mass $B^{-4}$ after summing over
coarse output positions, and each line convolution preserves this mass.
Jensen followed by summation proves physical $L^p$ contraction. Deleting
boundary data or restricting to an incomplete row family does not preserve
this statement without a separate column estimate.

\subsection{Finite moments and the aggregate cost}

For $1\le p<\infty$, define
\[
 c_{\Delta,p}=\sigma_\Delta^2
 \left(\sigma_\Delta^{-4}\sum_{y\in\Delta}a^4|F(y)|^p\right)^{1/p},
 \qquad v_{x\Delta}=|\Delta\cap\supp K_r(x,\cdot)|_a.
\]
H\"older and \eqref{eq:kernel} give
\begin{equation}
 r^2|(Q_2F)(x)|
 \le C\sum_\Delta r^{-2}v_{x\Delta}^{1-1/p}
       \sigma_\Delta^{4/p-2}c_{\Delta,p}.
 \label{eq:pbound}
\end{equation}
For $\sigma\le r$ the single-cell bound is again
$C(\sigma/r)^2c_{\Delta,p}$. For $\sigma\ge r$, the displayed bound is
at most $C(r/\sigma)^{2-4/p}c_{\Delta,p}$. The actual overlap in
\eqref{eq:pbound} is the more informative expression.

The single-cell gain cannot be used before summing the cells. If the whole
readset has assigned length $\sigma\ll r$ and only
$c_\Delta\le\eps$ is known, \eqref{eq:overlap} gives
\begin{equation}
 r^2|Q_2F|\le\eps(r/\sigma)^2.
 \label{eq:dense}
\end{equation}
There can be $O((r/\sigma)^4)$ cells, each with the single-cell factor
$(\sigma/r)^2$. A cell with $\sigma\asymp r$ has no scale-ratio
suppression at all. These are tests of an inference from native bounds,
not constructed original selected histories. In particular,
\eqref{eq:dense} is an upper-bound calculation, not a claimed lower bound
for a Wilson configuration satisfying all selectors.

There is also an aggregate bound at the critical four-dimensional energy
exponent. Cauchy--Schwarz and the kernel maximum give
\begin{equation}
 r^2|Q_2F(x)|
 \le r^2\left(\sum_y a^{-4}K_r(x,y)^2\right)^{1/2}
            \norm{F}_{L_a^2(\supp K_r)}
 \le C\norm{F}_{L_a^2(\supp K_r)}.
 \label{eq:L2aggregate}
\end{equation}
Thus a squared curvature bound of order $r^4g_k^2q^2$ on the actual readset
would suffice, through this estimate, to obtain old field size $r^2g_kq$.
The positive-filter argument supplies no such selected-body energy bound.
In particular an unweighted exponential reserve for a dimensionless energy
cannot silently be replaced by this scale-dependent assertion.

\subsection{Boundary and mixed junction rows remain separate}

The full $q_{\rm old}$ is mixed. It contains identity rows at literal bare
boundaries and rows whose plaquette completion uses adjacent levels.
Equation \eqref{eq:flatcurl} applies only to the complete uniform patches
specified above. To make this restriction exact, let $E_{\rm reg}^{(2)}$
insert those regular target plaquette rows into the full actual target
index set, and let $\mathscr C_\omega$ denote the actual linearized target
plaquette map, with each row's own mesh, frame, and prescribed adjacent-level
completion included. Define
\begin{equation}
 \mathcal B_\omega
 :=\mathscr C_\omega Dq_{\rm old}(0)
   -E_{\rm reg}^{(2)}Q_2d_a.
 \label{eq:boundary}
\end{equation}
This finite linear identity retains the remaining rows; it does not bound
them. No positive normalized smoothing kernel or $(\sigma/r)^2$ gain is
assigned to $\mathcal B_\omega$. A literal identity row reads its input
datum unchanged. Spatial interface commutators belong to this calculation.
For a fixed history its row labels are frozen, so differentiating this
linear map does not differentiate a real selector. Whole source derivatives
of the selected measure have their separate sharp-selector surface partners.

\section{Current responses and anchor sums in the linear channel}

\subsection{The current needs a separate input}

When both shifted plaquette inputs and their complete patches belong to the
same uniform hierarchy, a forward old difference gives
\begin{equation}
 r^3|D_{\rho,r}Q_2F(x)|
 \le r^2\{Q_2|F|(x+re_\rho)+Q_2|F|(x)\}.
 \label{eq:currentcrudebound}
\end{equation}
This costs $r^{-1}$ in physical units, although it has no microscopic
$a^{-1}$ loss. Thus a physical curvature bound $Cg_k$ gives only
$r^3|D_rQ_2F|\le Cr^2g_k$. The saved regular current substitution requires
$Cr^3g_k$. Calling the output coarse does not supply the missing $r$.

The telescoping line identity gives a sharper alternative with a different
input:
\begin{equation}
 D_{\rho,r}Q_2F=Q_2L_{\rho,B}D_{\rho,a}F
 =:Q_{3,\rho}G_\rho.
 \label{eq:Q3}
\end{equation}
A backward difference has the corresponding shifted positive filter, again
with every shifted row and its complete patch retained in that hierarchy.
$Q_3$ has mass one, maximum coefficient at most $C(a/r)^4$, and range
$Cr$. If
$j_{\Delta,\rho}=\sigma_\Delta^3\sup_\Delta|D_{\rho,a}F|$, then
\begin{equation}
 r^3|D_{\rho,r}Q_2F(x)|
 \le\sum_\Delta w^{(3)}_{x\Delta}
               (r/\sigma_\Delta)^3j_{\Delta,\rho}.
 \label{eq:currentoverlap}
\end{equation}
One small cell contributes $C(\sigma/r)j_{\Delta,\rho}$. The sum must
still be bounded. The flat Wilson current is a finite sum of backward
directional derivatives, with the corresponding component filters. A bound
only on their already summed fine divergence does not bound the separately
filtered terms unless the required cancellation is proved.

\subsection{All fixed parameter jets, with their actual anchor measures}

Let $Z^{(m)}(y;z_1,\ldots,z_m)$ be an input response kernel, on the actual
anchor index set with its actual counting measures. For a metric $d_*$ and
$\eta\ge0$,
define
\begin{equation}
 \mathcal A^{(m)}_\eta(Z;y)
 =\sum_{z_1,\ldots,z_m}
     \prod_{i=1}^m(1+d_*(y,z_i))^\eta
        |Z^{(m)}(y;z_1,\ldots,z_m)|.
\end{equation}
Positivity of the kernel and
$1+d_*(x,z)\le(1+d_*(x,y))(1+d_*(y,z))$ imply
\begin{equation}
 \mathcal A^{(m)}_\eta(Q_2Z;x)
 \le\sum_yK_r(x,y)(1+d_*(x,y))^{m\eta}
                    \mathcal A^{(m)}_\eta(Z;y).
 \label{eq:anchor}
\end{equation}
This follows by applying the triangle inequality entry by entry and
interchanging nonnegative sums. It therefore also holds for increasing finite
anchor sets whenever the right side is finite. The analogous $Q_3$ statement
proves the directional-current-jet version.

For old-mesh Euclidean distance $d_*(x,y)=|x-y|/r$, the extra incidence
factor is at most $(1+C)^{m\eta}$. For a mixed native metric one must keep
\begin{equation}
 \mathcal I_{m,\eta}(x)
 =\sum_yK_r(x,y)(1+d_*(x,y))^{m\eta}.
 \label{eq:incidence}
\end{equation}
A physical range $Cr$ need not have bounded mixed-metric diameter. These
formulas preserve an already controlled input anchor sum; they do not
manufacture one from a directional operator norm. They also do not prove
the slower-rate native sum in NM12. Ordinary constant annihilation is
preserved because $Q_2 1=1$, but the full quadratic-area/cubic cancellation
and root-phase completion still need the actual nonlinear source packet.

\section{Nonlinear remainders, references, and the old minimum}
\label{sec:nonlinear}

\subsection{An exact remainder ledger on each admitted uniform hierarchy}

Fix a common logarithmic branch for the actual successive raw maps, and
keep their actual retained representatives $A_t$. Define exactly
\begin{equation}
 \mathcal N_t(A_t)
 :=F_{a_{t+1}}(q_tA_t)-Q_{2,t}F_{a_t}(A_t).
\end{equation}
Iteration gives, on that uniform subhierarchy,
\begin{equation}
 F_r(q_{a\to r}A)
 =Q_{2,a\to r}F_a(A)
  +\sum_tQ_{2,a_{t+1}\to r}\mathcal N_t(A_t).
 \label{eq:nonlineartelescope}
\end{equation}
This is an exact identity in the chosen frames, not a bound. It keeps every
intermediate raw mean. For the full mixed determining map its boundary and
junction operations must also be inserted explicitly; the uniform formula
does not replace them. Differentiating \eqref{eq:nonlineartelescope} leaves
$Q_2D^mF$ and all derivatives of the literal nonlinear remainders, including
their intermediate field responses.

CB17--18 already gives a useful estimate for these remainders in a verified
regular common potential chart, with the stated $C^{1,\beta}$ and $H^{2+s}$
bounds and gauge straightening. Those hypotheses have not been verified
for \eqref{eq:actual} on all mixed walls. Applying that regular estimate
here without its chart would repeat the incoming-domain gap.

\subsection{The reference response is present at cubic order}

The explicit outer reference in \eqref{eq:rawmean} cancels in its flat
differential with the $P_i$ held fixed, but not in its whole jet. Set
$P_i=e^{X_i}$, $R=e^Y$, and
$\bar X=\sum_iw_iX_i$. Expansion of the literal formula through total
degree three gives
\begin{equation}
 \log M=\bar X+\frac1{12}
 \left([\bar X,[\bar X,Y]]
           -\sum_iw_i[X_i,[X_i,Y]]\right)+O(4).
 \label{eq:referencecubic}
\end{equation}
For completeness, the inner logarithm begins
\[
 X_i-Y-\tfrac12[X_i,Y]-\tfrac1{12}[X_i,[X_i,Y]]
       +\tfrac1{12}[Y,[Y,X_i]].
\]
Averaging it and expanding the final product with $e^Y$ cancels the
degree-two commutator and the double-$Y$ terms; the two terms displayed
in \eqref{eq:referencecubic} remain. In the actual raw map the $P_i$ themselves
are path and endpoint functions; their responses are also retained.
An exact free-algebra check verifies this coefficient independently by truncated
exponential and logarithm series. No original mixed-history realization is
claimed by this algebraic example. It shows why the flat positive kernel
cannot control all nonlinear current and anchor jets by itself.

\subsection{The old own-current jets once its chart is admitted}

In old physical units, write $a_{\rm old}$ for the unchanged fine mesh
relative to the old scale. On an admitted common small $C^1$ chart, the
actual Wilson current calculation LM38--41 gives
\begin{equation}
 J_{a_{\rm old}}(A)=M_{a_{\rm old}}A+N_{a_{\rm old}}(A),
 \qquad M_a=d_a^*d_a.
\end{equation}
The nonlinear map $N_a$ is analytic from that fixed $C^1$ ball to $C^0$,
uniformly in $a$. Indeed LM39 bounds $DN_a$ there, and $N_a(0)=0$;
radial integration bounds $N_a$. Cauchy on a fixed smaller ball then
bounds every fixed $D^nN_a$ uniformly in the microscopic mesh.

Assume the actual old minimum connection $A(v)$ stays in that smaller
$C^1$ ball with a positive margin independent of the microscopic mesh.
A labeled finite derivative set $I$ then satisfies the exact chain rule
\begin{align}
 D_I[J_a(A(v))]
 &=M_aD_IA(v)
   +\sum_{\pi\in\Part(I)}
     D^{|\pi|}N_a(A(v))[D_BA(v):B\in\pi],\\
 \norm{D_IJ_a(A(v))}_{C^0}
 &\le\norm{M_aD_IA(v)}_{C^0}
   +\sum_{\pi\in\Part(I)}C_{|\pi|}
                  \prod_{B\in\pi}\norm{D_BA(v)}_{C^1}.
 \label{eq:owncurrent}
\end{align}
Thus higher current jets can be transferred from the old minimum's
Maxwell and first-order response jets on that admitted chart. This statement
does not infer those jets for the physical connection associated with
$W=P_{\rm old}q_{\rm old}\Gamma^R$. In a logarithmic old chart its convention
is $W_b=\exp(a_{\rm old}A_b)$; the factor $a_{\rm old}$ is not reset.
The actual target margin, old branch, Maxwell responses, and weighted anchor
sums remain inputs. Constrained criticality does not set the physical
current or the retained boundary reactions to zero.

\section{A finite-jet path estimate with the mixed cost left explicit}

\begin{lemma}[Recorded-word jets]
Suppose a literal oriented word $\gamma$ has a verified relative representation
$U_b(v)=e^{a_bH_b(v)}\bar U_b$ with fixed unitary reference links.
For inverse-oriented links use their actual inverse/conjugated relative
exponents. Put
\[
 W_\gamma=U(\gamma)\bar U(\gamma)^{-1},\qquad
 \zeta_\gamma=\sum_{b\in\gamma}a_b\norm{H_b},\qquad
 \Xi_B(\gamma)=\sum_{b\in\gamma}a_b\norm{D_BH_b}.
\]
For every nonempty fixed labeled derivative set $I$,
\begin{align}
 \norm{W_\gamma-1}&\le e^{\zeta_\gamma}-1,\\
 \norm{D_IW_\gamma}
 &\le e^{\zeta_\gamma}
        \sum_{\pi\in\Part(I)}\prod_{B\in\pi}\Xi_B(\gamma).
 \label{eq:wordjets}
\end{align}
\end{lemma}

\begin{proof}
Moving the fixed reference links to the right expresses $W_\gamma$ as a
product of exponentials conjugated by fixed unitary prefixes. These
conjugations preserve all the norms in the statement. Telescoping proves
the zero-order bound. The Duhamel simplex formula and the product rule give,
for the product as a function of its list of exponents,
\[
 \norm{D^m\mathcal E[U_1,\ldots,U_m]}
 \le e^{\zeta_\gamma}
            \prod_{i=1}^m\sum_b a_b\norm{U_{i,b}}.
\]
The permutations cancel the simplex-volume factorial; no extra microscopic
bond count appears. Composing with $H_b(v)$ and grouping the chain-rule
terms by set partitions proves \eqref{eq:wordjets}.
\end{proof}

If the available response is only the native bound
$\sigma_b\norm{D_BH_b}\le h_{\sigma_b,B}$, then the actual cost is
\begin{equation}
 \Xi_B(\gamma)
 \le\sum_\sigma\frac{\ell_{\gamma,\sigma}}{\sigma}
                       h_{\sigma,B},
 \qquad
 \ell_{\gamma,\sigma}
 =\sum_{b\in\gamma:\sigma_b=\sigma}a_b.
 \label{eq:wordcost}
\end{equation}
The microscopic count has canceled, while the assigned-scale incidence
remains. This is a sufficient estimate to test, not a necessary lower
bound: cancellation may improve it. Its failure to be uniformly small
would not alone disprove old-target admission.

For actual anchor sums, fix one root $x_\gamma$. Let $\tau$ be the minimal
connecting-tree length to its anchors, and $R$ their maximum distance
from the root. The weight $e^{\mu\tau}(1+R)^p$ is submultiplicative
over a partition with this common root: union of the connecting trees gives
$\tau_{\rm union}\le\sum\tau_B$, and
$1+R_{\rm union}\le\prod(1+R_B)$. Hence Tonelli applies
\eqref{eq:wordjets} to the corresponding weighted sums of $D_BH_b$.
If available estimates are instead rooted separately at each input bond,
transporting them to $x_\gamma$ has its own distance cost. It must be
included; the common-root formula does not remove it.

Moving reference frames, compensating gauges, and endpoint words have
their derivatives in the actual product rule. They can be absorbed into
$H_b$ only after a common fixed-reference chart is proved. A Casimir or
other second-order differential operator must be expanded into its actual
ordered first-order derivatives, retaining all coefficient responses.
These field-parameter estimates are not weak integration-by-parts identities
for a selected measure. The latter still retain the complete inherited
score and sharp-boundary partners of~\cite[pp.~1200--1204]{W} and~\cite[\S5]{CE}.

\section{What first-entry geometry excludes, and what remains unclassified}

The resonance test $\sigma\asymp r$ does not itself produce a counterexample
to the requested lemma. AE10's pinned fine-star minimum is a genuine native
constrained Wilson minimum, but no original R-I selected history or stored
first-entry representative has been identified for it.

There is a real chronological exclusion for a literal surviving old live
cell. If that cell already lies in the unchanged complete source readset,
the mark has touched that live component at that earlier stage. Later live
residence cannot be counted as additional earned regular age. This applies
to the full record, not just to the source anchor.

For the no-new-failure survival sector, the original next-scale cover and
layer enlargement in R-II~\cite{rii}, p.\ 384, gives a quantitative version. If an
ancestor at stage $m$ stays at distance at most $D\ell_j$ from a fixed
point of the complete old record, and the enlargement contains a
$cM\ell_{n+1}$ neighborhood at step $n+1$, the point is reached by an
index bounded by
\begin{equation}
 \max\left\{m+1,\ j+\left\lceil\log_L\frac{D}{cM}\right\rceil+1\right\},
 \label{eq:firsttouch}
\end{equation}
up to the fixed initial shell convention. For $D=0$ the immediate first
case applies. The estimate follows by taking the first later scale with
$cM\ell_n\ge D\ell_j$.

If a completed carrier actually connects points separated by
$D\ge cL^{k-j}$ in birth units, its existing spare locality factor pays
any fixed saved age power:
\begin{equation}
 e^{-\kappa D}
 \le C_{s,\kappa,c}L^{-s(k-j)}e^{-\kappa D/2}.
\end{equation}
This is the boundedness of $u^se^{-\kappa cu/2}$ for $u=L^{k-j}$.
The complete carrier and its connection must already be present; one
cannot infer this payment for a remote dependence that has not been
localized and filed.

The missing step is a classification through the \emph{actual three joins}:
does every lower assigned-scale row read by \eqref{eq:actual} have such an
old live ancestor, or is it a newly introduced determining boundary row?
For the latter, its actual target and response must be bounded directly.
The identity-valued interface argument following R-II (1.53) applies to
its specified annulus. It does not say that every later joined boundary
has identity data. Small corrections in (1.61)--(1.65) are measured relative
to their joined background; smallness of that difference does not remove
the background's lower-scale curvature.

Nor is a finite halo available merely because the original source is local
in the background field. The actual transverse Wilson calculation in
\S115.3 gives nonzero remote target responses. The literal boundary trace
retains its defining nonlocal dependence. The complete localization and
anchor estimates have to carry that dependence.

\section{The distinguished interaction and reserve on the same body}

This check concerns the body to which an incoming estimate would be
attached. It supplies no estimate for the unknown insertion in
\eqref{eq:actual}. Keep the literal positive reference body $\mu_0$, all
its selectors, inherited variables, reference factors and absolute scalars.
On R-II's verified outer domain for the fixed selected record, pp.\ 379--380
compare the complex outer extension with this body, including the full
inherited interaction. This does not enlarge that domain to every additional
source-wall extension required by the outstanding marked theorem. The bound is
\begin{equation}
 T_zF=\int e^{\sigma_z}F\dd\mu_0,
 \qquad \norm{\sigma_z}_\infty
 \le\eps_k:=Cg_kR_k^6p_0(g_k)^2q_0(g_k)
 \quad(d=4).
 \label{eq:body}
\end{equation}
The fixed geometric constants are included in $C$; the logarithmic factors
remain displayed. For $\eps_k<\log2$, elementary positivity gives
\begin{equation}
 |T_zF|\le e^{\eps_k}\int|F|\dd\mu_0,\qquad
 |T_z1|\ge(2-e^{\eps_k})M_0,
 \quad M_0=\mu_0(1)>0.
 \label{eq:mass}
\end{equation}
The second inequality follows from
$|T_z1-M_0|\le(e^{\eps_k}-1)M_0$.
The genuine mass cancels in this same-body quotient only after an insertion
estimate is known. It remains in raw numerators and does not cancel an
independent global Wilson denominator.

The full inherited $V(X^-)$ is distinct from the clean proper-difference
vertex of IV10. Its small outer variation in \eqref{eq:body} does not
imply $|V(X^-)|\le Cg_k(1+E)^{e/2}$, nor does it prove a radial
physical-fluctuation disc for $V(B)-V(0)$. The latter disc, its radius,
and compatibility with every readout and wall need their own proof.
The constant term $V(0)$ cannot be discarded.

\begin{proposition}[Reserve transfer for an identified clean subinteraction]
Let $\mu_1$ be the complete positive selected endpoint body, with
$\tau,\theta>0$ and $N,E\ge0$. Suppose a
\emph{literal} real subinteraction $W=\sum_Yv_Y$ has already admitted
same-body bounds
\begin{equation}
 |v_Y|\le g_ka_Y\sqrt{E_Y},\quad
 \sum_Ya_Y\le c_1N,\quad
 \sum_Ya_YE_Y\le c_2E,
 \qquad a_Y,E_Y\ge0.
\end{equation}
If
\begin{equation}
 \int e^{\tau N+\theta E}\dd\mu_1\le B,
 \qquad c_1c_2g_k^2\le\tau\theta,
\end{equation}
then the exact interpolation
$\dd\mu_t=e^{(t-1)W}\dd\mu_1$, $0\le t\le1$, obeys
\begin{equation}
 \sup_{0\le t\le1}\int e^{\tau N/2+\theta E/2}\dd\mu_t\le B.
 \label{eq:reserve}
\end{equation}
\end{proposition}

\begin{proof}
Weighted Cauchy--Schwarz and Young give
\[
 |W|\le g_k\sqrt{c_1c_2NE}
 \le\tfrac\theta2E+\frac{c_1c_2g_k^2}{2\theta}N
 \le\tfrac\theta2E+\tfrac\tau2N.
\]
Multiplying by $e^{\tau N/2+\theta E/2}$ bounds the interpolated positive
integrand by the original reserve integrand. Every selector and scalar
factor stays in $\mu_1$.
\end{proof}

For complex walls the exact score relative to the corresponding real
interpolated body includes
$\sigma_z+(t-1)(W_z-W_0)$. It cannot be replaced by $\sigma_z$ alone.
Direct majorization works if the displayed bounds also hold for the actual
$W_z$. This proposition settles interpolation of an available reserve for
an identified interaction. It establishes neither the interaction bounds
nor the endpoint reserve in a general mixed history.

There is also a concrete reason to keep rarity unnormalized. The original
failed-test quadratic branch in R-II, pp.\ 381--383, selects an old
preparatory bond with $|B_j(b)|\ge A_1p_1(g_j)$. On any nonzero positive
subbody of that branch,
\begin{equation}
 \frac{\int e^{\theta|B_j(b)|^2}\dd\mu_\omega}
      {\int\dd\mu_\omega}
 \ge e^{\theta A_1^2p_1(g_j)^2}.
 \label{eq:failedtest}
\end{equation}
Any relative-bank constant on such a nonzero branch must therefore dominate
the right side of \eqref{eq:failedtest}. Raw mass smallness alone cannot
yield the required uniform relative bank. A categorical obstruction along
arbitrarily weak couplings would additionally require nonzero such branches
along that sequence; this existence has not been established here.
This does not contradict a current-only estimate or a global Wilson-normalized
bound. The old rare suppression must be kept through the actual history sum. The absolute
vacuum-convention comparison in S19--S23 and the full outer-history
covariance remain separate obligations.

\section{The concrete estimate still required}

The overlap calculation isolates quantities that could support the saved
regular substitution without demanding an unweighted norm for the entire
microscopic prepared field. On the regular uniform rows, define
\begin{align}
 \mathcal T_2(x)&=
   \sum_\Delta w_{x\Delta}\sigma_\Delta^{-2}c_\Delta,\\
 \mathcal T_{3,\rho}(x)&=
   \sum_\Delta w^{(3)}_{x\Delta}\sigma_\Delta^{-3}j_{\Delta,\rho}.
 \label{eq:targets}
\end{align}
The useful sufficient inputs would be bounds of order $Cg_kq$ for these
quantities and their saved anchor-weighted finite jets, together with the
literal nonlinear remainder, boundary/junction, and old-minimum response
bounds. They would retain the desired old sizes $r^2g_kq$ and $r^3g_kq$.
The raw triangle inequalities may be stronger than necessary; a proved
signed cancellation could be used instead. No such cancellation is assumed.

To close the incoming lemma, one must now prove on the original full
simultaneous domain, in the following order:
\begin{enumerate}
\item The actual old target, including every boundary and junction row,
lies inside the already saved old radius with a strict margin, on a common
branch of $P_{\rm old}$.
\item Its actual minimum has the finite Maxwell/current jets and anchor
sums in the saved norms. Equations \eqref{eq:anchor}, \eqref{eq:owncurrent},
and \eqref{eq:wordjets} specify response inputs; they do not prove them.
\item The same completed root and covariant quadratic/cubic packet has
its saved moment cancellation after this composition; its physical
quartic tail is evaluated inside its original radius.
\item The actual distinguished interaction has its required $g_k$ bound,
and one unused joint reserve holds on that same selected body, with all
absolute normalizers and inherited scores retained.
\end{enumerate}
Only after these steps can the raw connected estimate of~\cite[\S3]{CE} attach the
stored factor with the weights \eqref{eq:weights}. A globally normalized
$B/Z$ statement then additionally needs the genuine history sum and the
independent absolute denominator comparison.

\subsection{Primary input not used here}

The literal R-I~\cite{ri} preparation conditions, (1.64)--(1.69) at
$n=N$, selector (1.75), and the relevant cutoff in (1.82) are not
used here. R-II's range assertion does not supply
those missing clauses. The particular geometric question is whether the
lower assigned rows in the complete old readset have genuine previous
live ancestry or are artificial joined boundaries with a sufficiently
suppressed target. This needs the actual definitions, not an inferred
replacement domain. The full R-I definitions would make this check
possible; they would not themselves prove the marked estimate.

For the subordinate native anchor route, the 1985 paper~\cite{var85}
refers the mixed metric and summability to
1984 Propagators II~\cite{prop84}, pp.\ 231--234 and Lemma 2.1, for
slower-rate summability. That metric and lemma are not verified here. No
surrogate metric, or matching of differently named decay constants is
promoted to a verified native estimate. Even a verified native sum would
still need connection to the actual old source readout.

\subsection{Checks and final scope}

An exact-arithmetic finite calculation checks the noncommutative cubic
reference term from exponential/logarithm series, positive filter masses
and coefficient bounds, physical column identities, direct versus nested
filter composition, and the directional difference identity. These checks
support the finite algebra used above. They do not test the original
history selectors, nonlinear minimum, wall-domain admission, or saved
three-row bound.

\begin{center}
\begin{tabularx}{\linewidth}{@{}p{0.36\linewidth}X@{}}
\toprule
\textbf{Interface}&\textbf{Status}\\
\midrule
Uniform old-row flat averaging & Proved, including exact overlaps and linear
current/anchor transport.\\
Nonlinear and reference terms & Exact identities recorded; full mixed-domain
bounds remain unproved.\\
Old boundary and junction rows & Retained explicitly; no smoothing gain or
admission asserted.\\
Saved LM42 three-row norm & Unproved on the original simultaneous touched
mixed domain.\\
Same-body interaction and reserve & Exact transfer available for admitted
inputs; the general mixed inputs remain unverified.\\
Global $B/Z$ and outer history & Not supplied by a conditional component
comparison; remain open.\\
\bottomrule
\end{tabularx}
\end{center}

\newpage
\begin{thebibliography}{9}
\bibitem{W} Wilson-source estimates, volumes 1--10 (\texttt{lattice\_theory/source\_estimates} in this collection), cited by source page. Especially \S\S45.4, 97.2, 100.10--100.12,
107.10--107.13, 115.3, 123 (AE8), 141, 154 and 159 (pp.~329, 667, 692--697, 778--785, 845--847, 932--935, 1070--1081, 1173--1180, 1214--1218).
\bibitem{CE} \emph{Connected estimates for selected operations}, \texttt{lattice\_theory/block\_and\_transfer/conditional\_estimates/connected\_estimates\_for\_selected\_operations.tex} in this collection.
Raw selected-body connected estimates and the unresolved incoming
representative. The source, scalar convention, and limitations are
preserved here.
\bibitem{sf1} T. Balaban, \emph{Renormalization group approach to lattice
gauge field theories. I. Generation of effective actions in a small field
approximation and a coupling constant renormalization in four dimensions},
Commun. Math. Phys. \textbf{109} (1987), 249--301.
\href{https://doi.org/10.1007/BF01215223}{doi:10.1007/BF01215223}.
Raw-mean formulas on pp.\ 253--254.
\bibitem{ri} T. Balaban, \emph{Large field renormalization. I. The basic
step of the R operation}, Commun. Math. Phys. \textbf{122} (1989), 175--202.
\href{https://doi.org/10.1007/BF01257412}{doi:10.1007/BF01257412}.
Cited for identification only; its range and selector
clauses are not used here.
\bibitem{rii} T. Balaban, \emph{Large field renormalization. II.
Localization, exponentiation, and bounds for the R operation},
Commun. Math. Phys. \textbf{122} (1989), 355--392.
\href{https://doi.org/10.1007/BF01238433}{doi:10.1007/BF01238433}.
Particularly pp.\ 370--384.
\bibitem{var85} T. Balaban, \emph{The variational problem and background
fields in renormalization group method for lattice gauge theories},
Commun. Math. Phys. \textbf{102} (1985), 277--309.
\href{https://doi.org/10.1007/BF01229381}{doi:10.1007/BF01229381}.
Geometry and cited response estimates.
\bibitem{prop84} T. Balaban, \emph{Propagators and renormalization
transformations for lattice gauge theories. II}, Commun. Math. Phys.
\textbf{96} (1984), 223--250.
\href{https://doi.org/10.1007/BF01240221}{doi:10.1007/BF01240221}.
The metric and slow-rate summability target are not verified here.
\end{thebibliography}
\end{document}
